% DEPRECATED — superseded by latex/deliverable/. Terminology (e.g. "algorithmic amplification")
% is outdated; current framing uses "out-of-network diffusion".
% =============================================================================
%  BlueCCGL — Project Plan
%  Bluesky Contrastive Cascade Graph Learning
% =============================================================================
\documentclass[12pt,a4paper]{article}

% ---- Packages ---------------------------------------------------------------
\usepackage[utf8]{inputenc}
\usepackage[T1]{fontenc}
\usepackage{lmodern}
\usepackage[english]{babel}
\usepackage{geometry}
\geometry{margin=2.5cm}

\usepackage{amsmath,amssymb}
\usepackage{graphicx}
\usepackage{booktabs}          % professional tables
\usepackage{longtable}         % tables spanning pages
\usepackage{array}             % extended column definitions
\usepackage{enumitem}          % list customisation
\usepackage{xcolor}            % colour for warning flags
\usepackage{hyperref}
\hypersetup{
  colorlinks  = true,
  linkcolor   = blue!70!black,
  citecolor   = green!50!black,
  urlcolor    = blue!60!black,
}
\usepackage[
  backend=biber,
  style=apa,
  sorting=nyt,
  natbib=true
]{biblatex}
\addbibresource{references.bib}

% ---- Custom commands --------------------------------------------------------
\newcommand{\flagcite}[1]{%
  {\color{orange}\textsuperscript{⚠}}\,#1%
}

% ---- Title ------------------------------------------------------------------
\title{\textbf{BlueCCGL: Bluesky Contrastive Cascade Graph Learning}\\[6pt]
       \large Project Plan}
\author{}   % Add author name here
\date{}     % Add date or leave empty

% =============================================================================
\begin{document}
\maketitle

% =============================================================================
\section{Introduction}
% =============================================================================

Understanding how information spreads through social networks is a central
question in computational social science.  On centralised platforms such as
Twitter/X and Weibo, information cascades---the chains of reposts through
which a piece of content reaches new audiences---have been studied extensively
using supervised deep learning methods
\citep{chen2019cascn, cao2017deephawkes, yuan2024casflow}.  However, these
approaches share two critical limitations: they require ground-truth labels
(e.g., ``viral'' vs.\ ``non-viral''), and they are designed for platforms
where a single opaque recommendation algorithm governs content visibility.
Neither assumption holds on decentralised social networks.

Bluesky, built on the AT Protocol, represents a new class of social platform
where users can choose or disable algorithmic feeds entirely, where identity
is portable across servers (Personal Data Servers), and where the full
firehose of platform activity is publicly accessible in real time.  This
architecture creates a unique natural experiment: on the same platform, some
content spreads organically through follower networks while other content is
amplified by user-selected algorithmic feeds---and crucially, the structural
trace of both pathways is observable.  Yet, to date, no study has analysed
information cascade dynamics on Bluesky at a structural level.

This research proposes \textbf{BlueCCGL} (Bluesky Contrastive Cascade Graph
Learning), a framework for the unsupervised discovery of information cascade
archetypes on Bluesky.  BlueCCGL extends the Contrastive Cascade Graph
Learning (CCGL) framework of \citet{xu2022ccgl} with three novel
contributions:
\begin{enumerate}[label=\textbf{(\arabic*)}]
  \item \emph{Feed-aware graph augmentation strategies} that respect the
        semantic distinction between organic and algorithmically-amplified
        spread;
  \item \emph{Edge-type aware message passing} via Relational Graph
        Convolutional Networks \citep[R-GCN;][]{schlichtkrull2018rgcn} that
        explicitly encodes whether a cascade edge represents follower-based or
        non-follower-based propagation; and
  \item \emph{Hierarchical graph pooling} \citep{ying2018diffpool} that
        captures multi-scale cascade structure from individual user decisions
        to community-level patterns.
\end{enumerate}
By learning cascade representations through contrastive self-supervised
learning and clustering them with HDBSCAN \citep{mcinnes2017hdbscan},
BlueCCGL discovers natural cascade archetypes without requiring any labelled
data.

The expected results are significant for two reasons.  First, a label-free
approach to cascade archetype discovery is directly transferable to any
platform where labelled cascade data is unavailable---which is the norm
outside a handful of well-studied centralised platforms.  Second, the
discovered archetypes would provide the first empirical characterisation of
how information spreads differently through follower networks versus
algorithmic recommendation on a decentralised, algorithm-optional
platform---offering insights that are impossible to obtain from centralised
platforms where these two pathways cannot be disentangled.

% =============================================================================
\section{Related Work}
% =============================================================================

% -----------------------------------------------------------------------------
\subsection{Information Cascade Prediction}
% -----------------------------------------------------------------------------

Graph neural network approaches to cascade modelling have evolved rapidly.
\textbf{CasCN} \citep{chen2019cascn} was among the first to apply GNNs to
cascade prediction, modelling cascades as sequences of graph snapshots with
temporal attention.  \textbf{DeepHawkes} \citep{cao2017deephawkes} bridged
classical Hawkes processes with deep learning to provide interpretable
attention over user influence, though it assumes homogeneous user influence.
More recently, \textbf{CasFlow} \citep{yuan2024casflow} introduced a
variational approach with hierarchical pooling, achieving state-of-the-art
performance on Weibo and Twitter datasets.
\flagcite{\textbf{HierCas} \citep{wang2024hiercas}} extended this with
multi-resolution modelling capturing both user-level and community-level
dynamics, and \flagcite{\textbf{TempCas} \citep{liu2025tempcas}} applied
Transformer architectures with continuous-time attention.

\paragraph{Gap.}
All of the above methods are \emph{supervised}, requiring ground-truth
popularity or virality labels.  None account for heterogeneous edge types
(e.g., follower vs.\ non-follower propagation), and all were evaluated
exclusively on centralised platforms (Twitter, Weibo).  BlueCCGL addresses
all three gaps: it is unsupervised, edge-type aware, and targets a
decentralised platform.

% -----------------------------------------------------------------------------
\subsection{Contrastive and Self-Supervised Graph Learning}
% -----------------------------------------------------------------------------

\textbf{CCGL} \citep{xu2022ccgl} is the direct methodological predecessor of
this work.  It introduced contrastive self-supervised pre-training for cascade
graphs using sub-cascade sampling for augmentation and an InfoNCE
graph-level contrastive loss, demonstrating 10--20\% improvement over
supervised baselines.  \textbf{GraphCL} \citep{you2020graphcl} provided a
systematic study of graph augmentation strategies (node dropping, edge
perturbation, attribute masking, subgraph sampling) that informs augmentation
design.  \textbf{GRACE} \citep{zhu2020grace} introduced node-level
corruption-based contrastive learning.  For temporal graphs,
\flagcite{\textbf{TGC} \citep{zhang2023tgc}} proposed time-aware augmentation
strategies, and \flagcite{\textbf{CSTGCN} \citep{li2024cstgcn}} introduced
multi-scale temporal contrasting for evolving structural patterns.

\paragraph{Gap.}
CCGL's augmentations are generic (random sub-cascade sampling) and carry no
domain semantics.  No existing contrastive graph learning method incorporates
augmentations that are aware of the distinction between organic and
algorithmically-mediated spread.  BlueCCGL introduces feed-aware augmentation
strategies (follower-subgraph extraction, temporal prefix sampling,
bridge-node removal) designed to preserve this semantic structure.

% -----------------------------------------------------------------------------
\subsection{Heterogeneous GNNs and Organic vs.\ Algorithmic Spread}
% -----------------------------------------------------------------------------

Relational and heterogeneous GNN architectures provide the building blocks
for edge-type aware cascade encoding.  \textbf{R-GCN}
\citep{schlichtkrull2018rgcn} introduced relation-specific weight matrices
with basis decomposition for parameter efficiency.  \textbf{HAN}
\citep{wang2019han} proposed hierarchical attention across node-level and
semantic-level, while \textbf{HGT} \citep{hu2020hgt} applied
Transformer-style attention to heterogeneous graphs.  \textbf{RHGNN}
\citep{yu2023rhgnn} demonstrated the importance of modelling multiple
relationship types (follow, mention, repost) in social network analysis.

On the empirical side, \citet{huszar2022twitter} provided a landmark internal
Twitter study showing algorithmic amplification of political content---a gold
standard that requires platform-insider access.  \citet{goel2016structural}
defined the ``structural virality'' metric (Wiener index) to distinguish
broadcast from viral spread, and \citet{vosoughi2018spread} showed that true
and false news produce structurally different cascades.
\flagcite{\citet{bhadani2024amplification}} proposed structural signatures of
algorithmic amplification, and \citet{bandy2021auditing} developed methodology
for auditing algorithmic curation.

\paragraph{Gap.}
No prior work has applied heterogeneous GNN encoding to information cascades,
specifically encoding follower vs.\ non-follower edges as distinct relation
types.  Furthermore, studies of organic vs.\ algorithmic spread have relied on
either platform-insider data \citep{huszar2022twitter} or indirect proxies.
BlueCCGL exploits Bluesky's observable follow graph to directly label cascade
edges and uses R-GCN encoding to learn relation-aware representations.

% -----------------------------------------------------------------------------
\subsection{Decentralised Social Networks and Bluesky}
% -----------------------------------------------------------------------------

Research on decentralised social networks has focused primarily on the
Fediverse.  \citet{zignani2018mastodon} conducted the first large-scale
Mastodon measurement study, and \citet{raman2019challenges} analysed technical
challenges of federation protocols.  For the AT Protocol specifically,
\citet{kleppmann2024atprotocol} provided a technical comparison of
ActivityPub and AT Protocol, including Bluesky's Firehose architecture.
\flagcite{\citet{cerisara2025bluesky}} presented an early Bluesky ecosystem
analysis studying user behaviour and AT Protocol adoption.

Most notably, \textbf{\citet{failla2024bluesky}} released the first
large-scale Bluesky dataset, covering the complete post history of over 4M
users ($\sim$81\% of all registered accounts at the time), 235M posts, a
directed follower graph of $\sim$144M edges, and timestamped
repost/quote/reply interaction graphs---all publicly available on Zenodo
(DOI:~\href{https://doi.org/10.5281/zenodo.14669616}{10.5281/zenodo.14669616}).
Critically, the dataset also includes posts from 11 feed generators and user
feed-subscription data, providing ground-truth information about algorithmic
exposure that is unavailable from the Jetstream firehose alone.

\paragraph{Gap.}
While \citet{failla2024bluesky} provide the first comprehensive Bluesky
dataset with descriptive analysis, no study has performed cascade-level
\emph{structural} analysis on Bluesky or any AT Protocol platform---specifically,
no work has modelled repost cascades as graphs, assigned edge types based on
follower relationships, or applied graph representation learning to discover
cascade archetypes.  BlueCCGL is designed to fill this gap, leveraging the
\citet{failla2024bluesky} dataset as its primary data source.

% -----------------------------------------------------------------------------
\subsection{Hierarchical Graph Pooling and Clustering}
% -----------------------------------------------------------------------------

\textbf{DiffPool} \citep{ying2018diffpool} introduced differentiable soft
cluster assignment for hierarchical graph representation learning.
\textbf{SAGPool} \citep{lee2019sagpool} provided a more efficient
attention-based alternative, and \textbf{MinCutPool}
\citep{bianchi2020mincutpool} offered a spectral-clustering-grounded
approach.  For cascades specifically, \flagcite{\textbf{HiCas}
\citep{zhou2023hicas}} applied hierarchical pooling to capture multi-scale
propagation patterns.  For unsupervised pattern discovery, \textbf{HDBSCAN}
\citep{mcinnes2017hdbscan} provides density-based clustering that does not
require pre-specifying the number of clusters, while \textbf{DEC}
\citep{xie2016dec} and \textbf{DMoN} \citep{tsitsulin2023dmon} offer
end-to-end deep clustering alternatives.

\paragraph{Gap.}
Hierarchical pooling has not been combined with contrastive learning for
cascade graphs, nor has it been applied in an edge-type aware setting.
BlueCCGL integrates DiffPool-style hierarchical pooling into the contrastive
cascade learning pipeline, enabling multi-scale representations that capture
micro (user), meso (community), and macro (global cascade shape) patterns.

% =============================================================================
\section{Research Questions}
% =============================================================================

\paragraph{Main Research Question.}
\begin{quote}
\emph{Can unsupervised contrastive graph learning discover meaningful
information cascade archetypes on Bluesky that reflect the interplay between
organic follower-based spread and algorithmic amplification?}
\end{quote}

\noindent
This is decomposed into three sub-questions:

\paragraph{SQ1: Archetype Discovery.}
What natural cascade archetypes emerge when applying HDBSCAN clustering to
cascade embeddings learned via contrastive self-supervised learning, and how
do these archetypes differ in structural properties (depth, breadth, branching
factor, follower-edge ratio, temporal spread pattern)?

\paragraph{SQ2: Edge-Type Aware Encoding.}
How does encoding follower vs.\ non-follower edges via relational graph
convolution (R-GCN) affect the quality and discriminative power of learned
cascade representations, compared to homogeneous GNN baselines?

\paragraph{SQ3: Feed-Aware Augmentation.}
Do domain-specific augmentation strategies---follower-subgraph extraction,
temporal prefix sampling, and bridge-node removal---produce higher-quality
cascade representations than generic graph augmentations (random node/edge
dropping) as measured by clustering coherence and downstream task performance?

% =============================================================================
\section{Expected Outcomes}
% =============================================================================

\paragraph{Outcome~1---A taxonomy of cascade archetypes on Bluesky.}
We expect HDBSCAN clustering over the learned embeddings to reveal distinct
cascade archetypes characterised by interpretable structural properties.
Based on prior cascade literature \citep{goel2016structural,
vosoughi2018spread}, we anticipate at minimum a distinction between
\emph{broadcast-like} cascades (shallow, wide, dominated by non-follower
edges suggesting algorithmic amplification) and \emph{viral-like} cascades
(deep, narrow, dominated by follower edges suggesting organic diffusion),
with potentially additional intermediate archetypes (e.g., mixed-spread
cascades with early organic seeding followed by algorithmic amplification).

\paragraph{Outcome~2---Quantitative evidence for the contribution of each
extension.}
Through ablation studies, we expect to demonstrate that:
\begin{enumerate}[label=(\alph*)]
  \item edge-type aware R-GCN encoding improves representation quality over
        homogeneous GNN baselines (measured via silhouette score,
        Davies--Bouldin index, and Calinski--Harabasz index);
  \item feed-aware augmentations outperform generic random augmentations on
        the same metrics; and
  \item hierarchical pooling captures multi-scale cascade patterns that flat
        global pooling misses.
\end{enumerate}
If any extension does not yield improvement, this is itself a valuable
finding about the nature of cascade structure on Bluesky.

\paragraph{Outcome~3---An open-source implementation and reusable cascade
dataset.}
The project will produce:
\begin{enumerate}[label=(\alph*)]
  \item the BlueCCGL PyTorch Geometric codebase (cascade construction, model
        training, clustering); and
  \item a processed Bluesky cascade graph dataset derived from the
        \citet{failla2024bluesky} Zenodo dataset
        (DOI:~\href{https://doi.org/10.5281/zenodo.14669616}{10.5281/zenodo.14669616}),
        containing approximately 50,000 cascades (10--10,000 reposts each)
        with edge-type annotations (follower / non-follower / unknown)
        assigned by cross-referencing repost edges against the dataset's
        144M-edge follower graph.
\end{enumerate}

\paragraph{Outcome~4---Transfer learning evaluation on downstream tasks.}
We expect the pre-trained BlueCCGL encoder to provide useful representations
for downstream cascade tasks---specifically cascade size prediction and
cascade longevity prediction---outperforming both a random-initialisation
baseline and vanilla CCGL, thereby demonstrating the practical utility of the
learned representations beyond archetype discovery.

% =============================================================================
\section{Research Plan}
% =============================================================================

This is a solo Bachelor's thesis project spanning 16~weeks.  By leveraging
the pre-existing \citet{failla2024bluesky} Zenodo dataset instead of live
Jetstream collection, the data acquisition phase is compressed from 4~weeks
to 2~weeks, freeing time for more thorough experimentation and writing.
Table~\ref{tab:schedule} details the planned activities per week.

\begin{longtable}{
  @{}
  >{\centering\arraybackslash}p{0.8cm}   % Week
  >{\raggedright\arraybackslash}p{2.6cm}  % Phase
  >{\raggedright\arraybackslash}p{7.2cm}  % Activities
  >{\raggedright\arraybackslash}p{4.0cm}  % Deliverable
  @{}
}
\caption{Planned weekly schedule for the BlueCCGL project.}
\label{tab:schedule} \\
\toprule
\textbf{Week} & \textbf{Phase} & \textbf{Activities} & \textbf{Deliverable} \\
\midrule
\endfirsthead

\multicolumn{4}{c}{\tablename\ \thetable{} -- \emph{continued from previous page}} \\
\toprule
\textbf{Week} & \textbf{Phase} & \textbf{Activities} & \textbf{Deliverable} \\
\midrule
\endhead

\midrule
\multicolumn{4}{r}{\emph{Continued on next page}} \\
\endfoot

\bottomrule
\endlastfoot

1 & Data Acquisition &
  Download \citet{failla2024bluesky} Zenodo dataset ($\sim$22\,GB);
  decompress \texttt{followers.csv.gz}, \texttt{interactions.csv.gz},
  \texttt{graphs.tar.gz}, \texttt{user\_posts.tar.gz}; load follower graph
  and interaction edge lists into working format (Parquet/DuckDB) &
  Raw dataset loaded; storage pipeline validated \\
\addlinespace

2 & Data Acquisition &
  Data quality checks (timestamp consistency, deduplication); compute
  exploratory statistics (repost volumes, cascade size distribution,
  follower-graph density, edge-type ratios); identify candidate cascades
  meeting size thresholds (10--10{,}000 reposts) &
  Data quality report; exploratory statistics notebook \\
\addlinespace

3 & Cascade Construction &
  Implement cascade graph builder: reconstruct repost cascades from
  \texttt{interactions.csv.gz} and \texttt{graphs.tar.gz}; assign edge types
  (follower / non-follower / unknown) by cross-referencing cascade edges
  against \texttt{followers.csv.gz} &
  Cascade construction pipeline; edge-type labelled cascades \\
\addlinespace

4 & Cascade Construction &
  Convert cascade graphs to PyTorch Geometric dataset; perform detailed EDA
  (cascade depth/breadth distributions, follower-edge ratio distributions,
  temporal spread patterns); filter to final $\sim$50K cascade dataset &
  PyG dataset; EDA notebook with visualisations \\
\addlinespace

5 & Model Implementation &
  Implement base CCGL reproduction: GNN encoder, sub-cascade sampling
  augmentation, InfoNCE contrastive loss &
  Base CCGL training loop verified on cascade data \\
\addlinespace

6 & Model Implementation &
  Implement Extension~1: feed-aware augmentation strategies
  (follower-subgraph extraction, temporal prefix sampling, bridge-node
  removal) &
  Augmentation module with unit tests \\
\addlinespace

7 & Model Implementation &
  Implement Extension~2: edge-type aware R-GCN encoder with basis
  decomposition; integrate with contrastive pipeline &
  R-GCN encoder integrated into BlueCCGL \\
\addlinespace

8 & Model Implementation &
  Implement Extension~3: hierarchical DiffPool-based pooling; assemble full
  BlueCCGL architecture; end-to-end training verification &
  Complete BlueCCGL model; successful training run \\
\addlinespace

9 & Experiments &
  Contrastive pre-training of BlueCCGL; hyperparameter tuning (temperature,
  augmentation ratios, learning rate) via Weights \& Biases &
  Trained model; hyperparameter sweep results \\
\addlinespace

10 & Experiments &
  Extract cascade embeddings; apply HDBSCAN clustering; compute clustering
  metrics (silhouette, Davies--Bouldin, Calinski--Harabasz); generate UMAP
  visualisations &
  Archetype clusters; clustering quality report \\
\addlinespace

11 & Experiments &
  Ablation studies: remove edge-type encoding, replace feed-aware
  augmentations with generic augmentations, replace hierarchical pooling
  with flat pooling, compare to random initialisation &
  Ablation results table \\
\addlinespace

12 & Experiments &
  Downstream transfer-learning evaluation (cascade size prediction, longevity
  prediction); archetype interpretation (structural property analysis per
  cluster) &
  Downstream task results; archetype characterisation \\
\addlinespace

13 & Writing &
  Draft Introduction, Related Work, and Research Questions sections; compile
  bibliography &
  Thesis draft chapters 1--3 \\
\addlinespace

14 & Writing &
  Draft Methodology and Experimental Setup sections; create all figures and
  tables &
  Thesis draft chapters 4--5 \\
\addlinespace

15 & Writing &
  Draft Results, Discussion, and Conclusion; integrate all sections into full
  thesis &
  Complete thesis draft \\
\addlinespace

16 & Writing &
  Revisions based on supervisor feedback; final proofreading; prepare
  presentation &
  Final thesis; presentation slides \\

\end{longtable}

% =============================================================================
% Flagged Citations
% =============================================================================

\subsection*{Flagged Citations Requiring Verification}

The following citations were included based on the initial literature survey
but require independent verification before final submission:

\begin{itemize}
  \item \citet{wang2024hiercas} --- ``Hierarchical Information Propagation
        Prediction''
  \item \citet{liu2025tempcas} --- ``Temporal Cascade Transformers''
  \item \citet{li2024cstgcn} --- ``Contrastive Self-Supervised Temporal Graph
        Learning''
  \item \citet{cerisara2025bluesky} --- ``Bluesky Ecosystem Analysis''
  \item \citet{zhou2023hicas} --- ``Hierarchical Cascade Representation''
  \item \citet{bhadani2024amplification} --- ``Algorithmic Amplification of
        Low-Quality Content''
\end{itemize}

% =============================================================================
% Bibliography
% =============================================================================

\newpage
\printbibliography

\end{document}
